\subsection{基数的指数运算}

定义 4. 设 a 和 b 为基数。从 b 到 a 的映射所构成的集合的基数记为 \(\mathfrak{a}^{\mathfrak{b}}\) ，此处为记号上的滥用。

这里符号的滥用在于， \(a^{b}\) 已经表示从 b 到 a 的映射的图像集合（第二章，§5，第 3 号），而这个集合不一定是基数（习题 2）。从上下文中总能清楚地看出该符号 \(a^{b}\) 应取哪种含义.

命题 9. 设 X 和 Y 是两个集合，a 和 b 分别是它们的基数。则集合 \(X^{Y}\) 的基数为 \(a^{b}\) .

因为存在一个从 \(\mathbf{X}^{\mathrm{Y}}\) 到映射集合（从 \(\mathfrak{b}\) 到 \(\mathfrak{a}\) 中的映射所成的集合）上的双射 (第二章，§5，第2号，命题2，推论)。

命题10。设 \(\mathfrak{a}\) 与 \(\mathfrak{b}\) 是基数，且令 \(\mathbf{I}\) 是一个集合，使得 Card (I) = \(\mathfrak{b}\) 。如果 \(\mathfrak{a}_{\iota} = \mathfrak{a}\) （对于所有 \(\iota \in \mathbf{I}\) ），我们有 \(\mathfrak{a}^{\mathfrak{b}} = \underset{\iota \in \mathbf{I}}{\mathsf{P}}\mathfrak{a}_{\iota}\) .

这由一族集合的乘积作为函数图集合的定义得出（第二章，§5，第3号）。

推论1. 设 \(\mathfrak{a}\) 为一个基数，并令 \((\mathfrak{b}_t)_{t\in I}\) 为一个基数族。则

\[
\mathfrak{a}^{\sum_{i \in I} \mathfrak{b}_i} = \underset {i \in I} {\mathbf {P}} \mathfrak{a}^{\mathfrak{b}_i}.
\]

设 S 为诸集合 \(\mathfrak{b}_i\) 的和，并令 \(a_s = a\) （对于所有 \(s \in S\) ）。待证等式的两边随后都等于 \(\underset{s \in S}{\mathsf{P}} a_s\) ，这依据命题10以及乘积的结合律公式（第3号，命题5(b)）。

推论2。设 \((\mathfrak{a}_i)_{i\in I}\) 是一族基数，且设 \(\mathfrak{b}\) 是一个基数。则

\[
\left(\underset {i \in I} {\mathbf {P}} a _ {i}\right) ^ {b} = \underset {i \in I} {\mathbf {P}} a _ {i} ^ {b}.
\]

令 \(a_{i\beta} = a_i\) （对于每一对 \((i, \beta) \in I \times b\) ）。那么，根据乘积的结合性，我们有

\[
\left(\underset {i \in I} {\mathsf {P}} a _ {i}\right) ^ {\mathfrak {b}} = \underset {\beta \in \mathfrak {b}} {\mathsf {P}} \left(\underset {i \in I} {\mathsf {P}} a _ {i \beta}\right) = \underset {i \in I} {\mathsf {P}} \left(\underset {\beta \in \mathfrak {b}} {\mathsf {P}} a _ {i \beta}\right) = \underset {i \in I} {\mathsf {P}} a _ {i} ^ {\mathfrak {b}}.
\]

推论3。设 \(a, b, c\) 是基数。则 \(a^{bc} = (a^{b})^{c}\) .

令 \(\mathfrak{b}_{\gamma} = \mathfrak{b}\) （对于所有 \(\gamma \in \mathfrak{c}\) ）。则

\[
a ^ {b c} = a ^ {\sum_ {\gamma \in c} b _ {\gamma}} = \underset {\gamma \in c} {\mathbf {P}} a ^ {b _ {\gamma}} = (a ^ {b}) ^ {c}
\]

根据推论1。

命题11。设 \(\mathfrak{a}\) 是一个基数。则 \(\mathfrak{a}^0 = 1\) , \(\mathfrak{a}^1 = \mathfrak{a}\) , \(1^{\mathfrak{a}} = 1\) ；以及 \(0^{\mathfrak{a}} = 0\) （若 \(\mathfrak{a} \neq 0\) ）.

因为存在唯一的从 \(\varnothing\) 到任意给定集合中的映射（即图形为空集的映射）；由单个元素组成的集合到任意集合X的映射之集与X等势（第二章，§5，第3号）；存在唯一一个从任意集合到由单个元素组成的集合的映射；最后，不存在从非空集合到 \(\varnothing\) 的映射。

请特别注意， \(0^{0}=1\) .

命题12。设X为一个集合，并设a为其基数。则 \(\mathfrak{P}(X)\) （\(X\) 的所有子集所成的集合）的基数是 \(2^{\mathfrak{a}}\) 。

令 \(\alpha\) 与 \(\beta\) 为基数2的元素。对于X的每个子集Y，令 \(f_{\mathbf{Y}}\) 为 X 到 2 中的映射，定义为 \(f_{\mathbf{Y}}(x) = \alpha\) （若 \(x \in \mathbf{Y}\) ）且 \(f_{\mathbf{Y}}(x) = \beta\) （若 \(x \in \mathbf{X} - \mathbf{Y}\) ）. 设 \(u\) 为映射 \(\mathbf{Y} \to f_{\mathbf{Y}}\) （从 \(\mathfrak{P}(\mathbf{X})\) 到 \(2^{\mathbf{x}}\) 中的映射）。反之，对于每个将X映入2的映射 \(g\) ，我们让它对应于X的子集 \(-g^{1}(\alpha)\) ，并且令 \(v\) 为映射 \(g \to -g^{1}(\alpha)\) （从 \(2^{\mathbf{x}}\) 到 \(\mathfrak{P}(\mathbf{X})\) 中的映射）. 显然， \(u \circ v\) 与 \(v \circ u\) 分别是 \(2^{\mathbf{x}}\) 与 \(\mathfrak{P}(\mathbf{X})\) 的恒等映射；因此 \(u\) 与 \(v\) 是双射（第二章，§3，第8号，命题8，推论），因此 Card \((\mathfrak{P}(\mathbf{X})) = 2^{\mathbf{a}}\) .

